\section{Polynomial moments of stretched exponentials}
\label{sec:stretched_exponential_summability}

The shell and tail estimates in \cite{OpenAI2026AreaLaw},
\texttt{09-amplification.tex}, lines 141--159 and 187--212, use
summability of polynomial factors against stretched-exponential decay.
For $p\in\mathbb N$ and $c,\alpha>0$, write
\begin{align}
 M_p(c,\alpha)&=\sum_{n=0}^{\infty}(1+n)^p e^{-cn^\alpha}.
 \label{eq:stretched_moment}
\end{align}
The next theorem establishes convergence of this literal series.
All the estimates in this section hold for every $\alpha>0$.
% Source revision: adc7f1241b42e322a6451854ab7e4b4c146bf78a.
% Checked module: QICLean/Analysis/StretchedExponentialSummability.lean.

\begin{theorem}[Summability of every polynomial moment]
\label{thm:stretched_moment_summable}
\lean{Real.summable_nat_pow_mul_exp_neg_mul_rpow}
\leanok
For every $p\in\mathbb N$, $c>0$, and $\alpha>0$, the series
in~(\ref{eq:stretched_moment}) is summable.
\end{theorem}
\begin{proof}
\leanok
The exponential $e^{-ct}$ decays faster than the power
$t^{(-p-2)/\alpha}$ as $t\to\infty$.  Substitute $t=n^\alpha$;
since $\alpha>0$, this tends to infinity and gives
$e^{-cn^\alpha}\leq n^{-p-2}$ for all sufficiently large $n$.
For these $n$, also require $n\geq1$.  Then
\begin{align}
 (1+n)^p e^{-cn^\alpha}
 &\leq (2n)^p n^{-p-2}=2^p n^{-2}.
 \label{eq:stretched_moment_comparison}
\end{align}
The summands are nonnegative, and $\sum_{n\geq1}n^{-2}$ converges.
Comparison outside a finite initial segment proves the claim.
\end{proof}

\begin{theorem}[Lower bound on the full moment]
\label{thm:stretched_moment_one_le}
\lean{Real.one_le_tsum_nat_pow_mul_exp_neg_mul_rpow}
\leanok
For every $p\in\mathbb N$, $c>0$, and $\alpha>0$, one has
$1\leq M_p(c,\alpha)$.
\end{theorem}
\begin{proof}
\leanok
\uses{thm:stretched_moment_summable}
The series is summable by Theorem~\ref{thm:stretched_moment_summable}
and has nonnegative summands.  Its $n=0$ term is exactly one,
because $0^\alpha=0$ for $\alpha>0$.
\end{proof}

\begin{theorem}[Tail estimate at a real cutoff]
\label{thm:stretched_moment_tail}
\lean{Real.tsum_nat_pow_mul_exp_neg_mul_rpow_tail_le}
\leanok
Let $p\in\mathbb N$, $c>0$, $\alpha>0$, and $r\in\mathbb R$
with $r\geq0$.  Then
\begin{align}
 \label{eq:stretched_moment_tail}
 \sum_{\substack{n\in\mathbb N\\r\leq n}}(1+n)^p e^{-cn^\alpha}
 &\leq
 \left(\sum_{n=0}^{\infty}(1+n)^p e^{-(c/2)n^\alpha}\right)
 e^{-(c/2)r^\alpha}.
\end{align}
The constant multiplying the decay at $r$ is the full moment
$M_p(c/2,\alpha)$, which is finite and at least one by
Theorems~\ref{thm:stretched_moment_summable}
and~\ref{thm:stretched_moment_one_le}.
The cutoff need not be an integer, and the boundary $n=r$ is included
whenever $r$ is a natural number.
\end{theorem}
\begin{proof}
\leanok
\uses{thm:stretched_moment_summable}
Both the full-rate and half-rate series are summable by
Theorem~\ref{thm:stretched_moment_summable}.  For every natural
number $n\geq r$, monotonicity of nonnegative real powers gives
$r^\alpha\leq n^\alpha$.  Splitting the exponential rate yields
\begin{align}
 \begin{aligned}
 (1+n)^p e^{-cn^\alpha}
 &=(1+n)^p e^{-(c/2)n^\alpha}e^{-(c/2)n^\alpha}\\
 &\leq (1+n)^p e^{-(c/2)n^\alpha}e^{-(c/2)r^\alpha}.
 \end{aligned}
 \label{eq:stretched_tail_comparison}
\end{align}
Sum over the natural numbers with $r\leq n$, enlarge to all natural
numbers using nonnegativity, and take the constant
$e^{-(c/2)r^\alpha}$ outside the sum.
\end{proof}
